\documentclass[11pt]{article}
\usepackage{amsmath,amssymb,amsthm}
\usepackage[margin=1in]{geometry}
\usepackage[hidelinks]{hyperref}
\newtheorem{lemma}{Lemma}\newtheorem{proposition}{Proposition}
\begin{document}
\setlength{\emergencystretch}{1em}
\title{The general quadratic-gap shift and a limitation of quadrant Pick positivity}
\author{}\date{}\maketitle
We give an exact recurrence for every lower exponent
$h$ at gap two.  It also disproves a proposed sufficient condition:
the actual shift quotient is not Pick on the right half-plane for any
$h\ge6$.  This is not a counterexample to consecutive coprimality or SFC.

Fix $n\ge1$ and $\gamma>0$, set $M=\gamma+2n$, and put
\[
T_{n,\gamma}(t)=\sum_{j=0}^n\binom nj(\gamma)_{2j}t^j,\quad
R(z)=z^{2n}T_{n,\gamma}(-z^{-2}),\quad E=M-zD_z,\quad u=z^2.
\]
Write $K_j=(E)_jR$.  The lower equation is
\[
EK_1=K_2-K_1,\qquad EK_2=MK_2+uK_1-\gamma uR.       \tag{1}
\]
Consequently the action of $E$ on a coefficient vector $(a,b,c)$
in the ordered basis $(R,K_1,K_2)$ is
\[
E(a,b,c)=(-2ua'-\gamma uc,\ a-b-2ub'+uc,\ b+Mc-2uc').
                                                               \tag{2}
\]
Derivatives here are formal derivatives in $u$.

\begin{lemma}[An exact exponent recurrence]\label{lem:shift-recurrence}
For every $h\ge0$ define polynomials $r_h,W_h,U_h$ by
\[
[u(E)_h-(E)_{h+2}]R=r_hR-W_hK_1-U_hK_2.
\]
One may choose them with the initial values $(r_0,W_0,U_0)=(u,0,1)$
and, for every $h\ge0$,
\[
\begin{split}
r_{h+1}&=(h+2)r_h-2ur_h'+(M-2n)uU_h,\\
W_{h+1}&=-r_h+(h+1)W_h-2uW_h'+uU_h,\\
U_{h+1}&=W_h+(M+h+2)U_h-2uU_h'.
\end{split} \tag{3}
\]
Their coefficient ring is $\mathbb Z[M,n,u]$.
For $h\ge1$ and $\gamma>0$, every common root $t$ of
$T_{n,\gamma}$ and $T_{n+1,\gamma+h}$ gives $u=-1/t\ne0$ and
\[
U_h(u)C_i(u)+W_h(u)B_i(u)=0,\quad i=0,1,             \tag{4}
\]
where $B=B_{n,\gamma}$, $C=B_{n+1,\gamma}$ and $B_i,C_i$
are their even--odd sections.
\end{lemma}
\begin{proof}
Let $Q_h=u(E)_h-(E)_{h+2}$.
Since $(E+h+2)u=u(E+h)$, one has
$Q_{h+1}=(E+h+2)Q_h$.  Equation (2) now proves (3), including
all parameter shifts and signs.
The exact Gamma-shift/Pascal identity is
\[
(\gamma)_h z^{2(n+1)}T_{n+1,\gamma+h}(-z^{-2})=Q_hR.
\]
At a common root $R=0$, this gives $W_hK_1+U_hK_2=0$.
The full concomitant of
\href{LOWER_TWO_PRIME_GAP_SFC.tex}
{Theorem \texttt{thm:mixed}, proof equation(3), at $d=2$},
gives $B(\pm z)K_2-C(\pm z)K_1=0$.
The full cubic jet at $z\ne0$ is nonzero; hence at $R=0$,
$(K_1,K_2)\ne(0,0)$.
Its determinant with the upper relation gives
$U_hC(\pm z)+W_hB(\pm z)=0$.  Even--odd inversion proves (4).
No division by $U_h$, $B$, or $R'$ was made.
\end{proof}

\begin{lemma}[A subtraction-free coefficient recurrence]\label{lem:positive-coefficients}
Put $s=2(n+1)$ and $V_h=uU_h-r_h$.  Then
\[
\begin{split}
U_{h+1}&=(M+h+2-2uD_u)U_h+W_h,\\
W_{h+1}&=(h+1-2uD_u)W_h+V_h,\\
V_{h+1}&=(h+2-2uD_u)V_h+uW_h+suU_h.
\end{split} \tag{7}
\]
For all $h\ge0$, $U_h,W_h,V_h$ belong to
$\mathbb Z_{\ge0}[M,s,u]$.
More precisely, $U_0=1$, $W_0=V_0=0$, and for $h\ge1$
\[
\deg_u U_h\le\left\lfloor\frac{h-1}{2}\right\rfloor,\quad
\deg_u W_h\le\left\lfloor\frac h2\right\rfloor,\quad
\deg_u V_h\le\left\lfloor\frac{h+1}{2}\right\rfloor,                 \tag{8}
\]
with the zero polynomial allowed in these bounds.
\end{lemma}
\begin{proof}
Substitute $r_h=uU_h-V_h$ into (3).  In the third equation,
differentiating $uU_h$ supplies the extra $2uU_h$;
the coefficient becomes $M-\gamma+2=s$, proving (7).
For $h=0$ the first step is
$(U_1,W_1,V_1)=(M+2,0,su)$, giving (8).
Induction on (7) proves the degree bounds.
At degree $u^j$ the diagonal multipliers are, respectively,
$M+h+2-2j$, $h+1-2j$, and $h+2-2j$.
Their integer parts are nonnegative wherever the corresponding
previous-row coefficient is permitted by (8).
All other terms in (7) are additions or multiplication by $u$ or $su$.
Thus coefficient nonnegativity follows by induction.

Equivalently, use a layered directed graph with vertices $(h,j,\tau)$,
$\tau\in\{U,W,V\}$.  The diagonal edges have the three displayed
weights; the other edges are $W\to U$, $V\to W$, $W\to V$
with the last one raising $j$ by one, and $U\to V$ raising $j$
with weight $s$.  Start at $(0,0,U)$.
The path sums are exactly the coefficients in (7).
This is a positive path-sum construction; no planarity or
total-positivity assertion is being made.
\end{proof}

For example, with $s=2(n+1)$,
\[
\begin{array}{c|c|c}
h&U_h&W_h\\ \hline
2&(M+2)_2&su\\
3&(M+2)_3+su&s(M+4)u\\
4&(M+2)_4+s(2M+7)u&su((M+4)_2+u)\\
5&(M+2)_5+3s(M+4)^2u+su^2&
su((M+4)_3+(M+2n+8)u).
\end{array}
\]
For the physical phase $X^h(A+BX^2)$ one substitutes
$\gamma=hn+\beta$, thus $M=(h+2)n+\beta$, $\beta>0$.
When $n\ge2$ and $\gamma>2n-2$, the exact quotient $C_1/B_1$ is strictly
Pick at nonreal points, by the general Pick lemma in the
\href{LOWER_FIVE_QUADRATIC_SECTOR_SFC.tex}
{lower-five sector note}.  Therefore (4), if $U_h\ne0$,
forces $\Im(W_h/U_h)/\Im u<0$.  It is tempting to exclude this
through right-half-plane Pick positivity; the following proposition
shows that this particular sufficient condition is false.

\begin{proposition}[An obstruction for every $h\ge6$]\label{prop:quadrant-obstruction}
Fix any integers $h\ge6$ and $n\ge1$.  For all sufficiently large
positive $\beta$ there is a nonreal $u$ with $\Re u>0$ such that,
at $M=(h+2)n+\beta$,
\[
U_h(u)\ne0,\qquad
\frac{\Im(W_h/U_h)(u)}{\Im u}<0.                     \tag{5}
\]
Thus full quadrant Pick positivity fails inside the physical
positive-Gamma family at every lower exponent $h\ge6$.
\end{proposition}
\begin{proof}
Let $\varepsilon=1/M$ and set
\[
\mathsf U_h=\varepsilon^hU_h(M^2v),\quad
\mathsf W_h=\varepsilon^hW_h(M^2v),\quad
\mathsf r_h=\varepsilon^{h+2}r_h(M^2v).
\]
For fixed $h,n$ these are polynomials in $\varepsilon,v$:
the only possible negative power in the transformed recurrence is
\[
\mathsf W_{h+1}
 =\varepsilon^{-1}(v\mathsf U_h-\mathsf r_h)
   +\varepsilon((h+1)\mathsf W_h-2v\partial_v\mathsf W_h).
\]
Inductively $\mathsf U_h(0,v)=1$ and $\mathsf r_h(0,v)=v$, so
the numerator is divisible by $\varepsilon$.
The other transformed recurrences are
\[
\begin{split}
\mathsf U_{h+1}
 &=\mathsf U_h+\varepsilon((h+2)\mathsf U_h
                -2v\partial_v\mathsf U_h+\mathsf W_h),\\
\mathsf r_{h+1}
 &=(1-2n\varepsilon)v\mathsf U_h+
        \varepsilon((h+2)\mathsf r_h-2v\partial_v\mathsf r_h).
\end{split}
\]
Write $a_h=[\varepsilon]\mathsf U_h$,
$b_h=[\varepsilon]\mathsf r_h$,
$w_h=\mathsf W_h(0,v)$.  The initial values vanish and
\[
a_{h+1}=a_h+h+2+w_h,\quad
b_{h+1}=va_h+(h-2n)v,\quad w_{h+1}=va_h-b_h.
\]
Hence $w_0=w_1=0$ and
$w_{h+1}=sv+vw_{h-1}$ for $h\ge1$.  In particular
\[
\lim_{M\to+\infty}\frac{W_h(M^2v)}{U_h(M^2v)}
 =s\sum_{j=1}^{\lfloor h/2\rfloor}v^j.              \tag{6}
\]
The denominator is nonzero for sufficiently large $M$, since
$\mathsf U_h\to1$.

Put $m=\lfloor h/2\rfloor\ge3$, $\theta=5\pi/(4m)$, and
$v=4e^{i\theta}$.  Then $0<\theta<\pi/2$ and
\[
\Im\sum_{j=1}^mv^j
 \le-\frac{4^m}{\sqrt2}+\sum_{j=1}^{m-1}4^j
 <-\frac{4^m}{\sqrt2}+\frac{4^m}{3}<0.
\]
By (6), $u=M^2v$ has the properties in (5) for all sufficiently
large $M$.  Such $M$ corresponds to $\beta>0$.
\end{proof}

The offending $u$ need not be a lower moment root.  This obstruction
leaves open a restriction of the quotient inequality to actual roots,
a sharper coupled condition using both cyclic rows, and other
all-exponent SFC arguments.  It also leaves open every unrestricted
binomial assertion. The exact controls accompanying this note check formulas
and a literal admissible quotient counterexample. The written induction
and limiting argument supply the unrestricted exponent quantifiers.
Internal execution is distinct from independent mathematical replay,
formal verification, external review and priority.
\end{document}
